\scanpage{43}
% archive sheet number: 42
\mspage{79}
\hfill n\textsuperscript{o} 382

$|\alpha|\leq k$ (cf. n\textsuperscript{o} 2.5.3 et 3.5.2) ; si $r\ne\omega$, ces valeurs peuvent être choisies arbitrairement.

L'ensemble $J^k(K^p,K^q)$ s'identifie canoniquement à $K^p\times K^q\times L^k_{p,q}$, où l'on a posé $L^k_{p,q}=J^k_{\unsure{o|0}}(K^p,K^q)_o$. L'ensemble $L^k_{p,q}$ admet une structure d'espace vectoriel sur $K$, caractérisée par
\[
j^k_o(f)+j^k_o(g)=j^k_o(f+g)\;,\quad\alpha.j^k_o(f)=j^k_o(\alpha.f)\qquad(\alpha\in K)\;,
\]
$f$ et $g$ étant des applications dans $K^q$ d'un voisinage ouvert de l'origine de $K^p$, de classe $C^r$, envoyant l'origine sur l'origine.

Supposons $k$ fini. Alors $L^k_{p,q}$ s'identifie canoniquement à $K^{q,c(k,p)}$, avec $c(k,p)=\binom{k+p}{p}$, sa base canonique étant formée des parties principales d'ordre $k$ des applications polynomiales $p_{j,\handwritten{\beta}}:x\longrightarrow(f^1(x),\ldots,f^q(x))$, où $f^i(x)=\handwritten{\delta}^i_j.x^{\handwritten{\beta}}$ ($j=1,\ldots,q$ ; $|\handwritten{\beta}|\leq k$). Cette identification est compatible avec la structure vectorielle définie précédemment. Les espaces $L^k_{p,q}$ et $J^k(K^p,K^q)$ sont ainsi munis canoniquement d'une structure de variété analytique.

L'ensemble $L^k_p$ des jets inversibles d'ordre $k$ de $K^p$ dans $K^p$, de source et but l'origine, muni de la loi de composition des jets, est une variété de groupe de classe $C^\omega$ ; le groupe $L^k_p\times L^k_q$ opère analytiquement à gauche sur $L^k_{p,q}$ par la composition des jets :
\[
(a,b).c=a\circ c\circ b^{-1}\qquad (a\in L^k_p,\ b\in L^k_q,\ c\in L^k_{p,q})\;.
\]

Soit $k'\leq k$. L'application naturelle $L^k_{p,q}\longrightarrow L^{k'}_{p,q}$ (cf. 7.1.4) est surjective, linéaire. Son noyau est l'ensemble des jets d'ordre $k$ de $K^p$ dans $K^q$ qui ont un contact d'ordre $\geq k'$ avec l'application nulle. L'application naturelle $L^k_p\longrightarrow L^{k'}_p$ est un morphisme surjectif de variétés de groupe de classe $C^\omega$.

7.1.6. Soient $X$ et $Y$ des variétés de classe $C^r$, $k$ un entier fini $\leq r$ et $k'$ un entier $\leq k$. Si $(U,\varphi,p)$ et $(V,\handwritten{\psi},q)$ sont des cartes de $X$ et $Y$, on a des isomorphismes
